\documentclass[supercite]{Experimental_Report}

\title{~~~~~~计算机系统基础~~~~~~}
\author{你的姓名}
\school{计算机科学与技术学院}
\classnum{你的专业班级}
\stunum{你的学号}
\instructor{你的指导教师}
\date{2026年X月X日}

\usepackage{algorithm, multirow}
\usepackage{algpseudocode}
\usepackage{amsmath}
\usepackage{amsthm}
\usepackage{bm}
\usepackage{float}
\usepackage{framed}
\usepackage{mathtools}
\usepackage{subcaption}
\usepackage{tikz}
\usepackage{tikzscale}
\usepackage{pgfplots}
\usepackage{listings}
\usepackage{xcolor}

\pgfplotsset{compat=1.16}

\lstdefinelanguage{Shell}{
  keywords={cd, ls, pwd, cat, grep, make, gcc, gdb, readelf, objdump, hexdump, python, echo, rm},
  sensitive=true,
  comment=[l]{\#},
  morestring=[b]',
  morestring=[b]"
}

\lstdefinestyle{reportcode}{
  basicstyle=\ttfamily\footnotesize,
  breaklines=true,
  numbers=left,
  numberstyle=\tiny,
  frame=single,
  captionpos=b,
  keywordstyle=\color{blue}\bfseries,
  stringstyle=\color{red},
  commentstyle=\color{gray}\itshape,
  showstringspaces=false,
  tabsize=2,
  language=C
}

\lstset{style=reportcode}

\newcommand{\cfig}[3]{
  \begin{figure}[htb]
    \centering
    \includegraphics[width=#2\textwidth]{images/#1.tikz}
    \caption{#3}
    \label{fig:#1}
  \end{figure}
}

\newcommand{\sfig}[3]{
  \begin{subfigure}[b]{#2\textwidth}
    \includegraphics[width=\textwidth]{images/#1.tikz}
    \caption{#3}
    \label{fig:#1}
  \end{subfigure}
}

\newcommand{\xfig}[3]{
  \begin{figure}[htb]
    \centering
    #3
    \caption{#2}
    \label{fig:#1}
  \end{figure}
}

\newcommand{\rfig}[1]{\autoref{fig:#1}}
\newcommand{\ralg}[1]{\autoref{alg:#1}}
\newcommand{\rthm}[1]{\autoref{thm:#1}}
\newcommand{\rlem}[1]{\autoref{lem:#1}}
\newcommand{\reqn}[1]{\autoref{eqn:#1}}
\newcommand{\rtbl}[1]{\autoref{tbl:#1}}

\algnewcommand\Null{\textsc{null }}
\algnewcommand\algorithmicinput{\textbf{Input:}}
\algnewcommand\Input{\item[\algorithmicinput]}
\algnewcommand\algorithmicoutput{\textbf{Output:}}
\algnewcommand\Output{\item[\algorithmicoutput]}
\algnewcommand\algorithmicbreak{\textbf{break}}
\algnewcommand\Break{\algorithmicbreak}
\algnewcommand\algorithmiccontinue{\textbf{continue}}
\algnewcommand\Continue{\algorithmiccontinue}
\algnewcommand{\LeftCom}[1]{\State $\triangleright$ #1}

\newtheorem{thm}{定理}[section]
\newtheorem{lem}{引理}[section]

\colorlet{shadecolor}{black!15}

\theoremstyle{definition}
\newtheorem{alg}{算法}[section]

\def\thmautorefname~#1\null{定理~#1~\null}
\def\lemautorefname~#1\null{引理~#1~\null}
\def\algautorefname~#1\null{算法~#1~\null}

\newcommand{\templatehint}[1]{{\color{gray}\textit{#1}}}

\begin{document}

% 仅用于编译和预览插图/表格（如 TikZ 等画图代码）
\usetikzlibrary{shapes.geometric, arrows.meta, positioning}

\begin{figure}[H]
  \centering
  \begin{tikzpicture}[
    node distance=1.5cm and 2cm,
    box/.style={rectangle, draw, rounded corners, minimum width=2.5cm, minimum height=1cm, align=center, fill=black!5},
    arrow/.style={-Latex, thick}
  ]
  
  \node[box] (data) {读取数据};
  \node[box, right=of data] (encode) {标签编码};
  \node[box, right=of encode] (split) {10 折划分};
  \node[box, below=of split] (train) {训练弱分类器};
  \node[box, left=of train] (update) {AdaBoost \\ 更新权重};
  \node[box, left=of update] (predict) {集成预测};
  \node[box, below=of predict] (output) {输出结果文件};
  
  \draw[arrow] (data) -- (encode);
  \draw[arrow] (encode) -- (split);
  \draw[arrow] (split) -- (train);
  \draw[arrow] (train) -- (update);
  \draw[arrow] (update.north) to[bend left] node[midway, above, font=\footnotesize] {未达T轮} (train.north);
  \draw[arrow] (update) -- node[midway, above, font=\footnotesize] {完成T轮} (predict);
  \draw[arrow] (predict) -- (output);
  
  \end{tikzpicture}
  \caption{AdaBoost 实验整体流程图}
  \label{fig:adaboost-flow}
\end{figure}

\end{document}
